% Options for packages loaded elsewhere
% Options for packages loaded elsewhere
\PassOptionsToPackage{unicode}{hyperref}
\PassOptionsToPackage{hyphens}{url}
\PassOptionsToPackage{dvipsnames,svgnames,x11names}{xcolor}
%
\documentclass[
  french,
  10pt,
  letterpaper,
  DIV=11,
  numbers=noendperiod]{scrartcl}
\usepackage{xcolor}
\usepackage[top=1.5cm,bottom=2cm,left=1.8cm,right=1.8cm]{geometry}
\usepackage{amsmath,amssymb}
\setcounter{secnumdepth}{-\maxdimen} % remove section numbering
\usepackage{iftex}
\ifPDFTeX
  \usepackage[T1]{fontenc}
  \usepackage[utf8]{inputenc}
  \usepackage{textcomp} % provide euro and other symbols
\else % if luatex or xetex
  \usepackage{unicode-math} % this also loads fontspec
  \defaultfontfeatures{Scale=MatchLowercase}
  \defaultfontfeatures[\rmfamily]{Ligatures=TeX,Scale=1}
\fi
\usepackage{lmodern}
\ifPDFTeX\else
  % xetex/luatex font selection
  \setmainfont[]{TeX Gyre Pagella}
  \setsansfont[]{TeX Gyre Heros}
\fi
% Use upquote if available, for straight quotes in verbatim environments
\IfFileExists{upquote.sty}{\usepackage{upquote}}{}
\IfFileExists{microtype.sty}{% use microtype if available
  \usepackage[]{microtype}
  \UseMicrotypeSet[protrusion]{basicmath} % disable protrusion for tt fonts
}{}
\makeatletter
\@ifundefined{KOMAClassName}{% if non-KOMA class
  \IfFileExists{parskip.sty}{%
    \usepackage{parskip}
  }{% else
    \setlength{\parindent}{0pt}
    \setlength{\parskip}{6pt plus 2pt minus 1pt}}
}{% if KOMA class
  \KOMAoptions{parskip=half}}
\makeatother
% Make \paragraph and \subparagraph free-standing
\makeatletter
\ifx\paragraph\undefined\else
  \let\oldparagraph\paragraph
  \renewcommand{\paragraph}{
    \@ifstar
      \xxxParagraphStar
      \xxxParagraphNoStar
  }
  \newcommand{\xxxParagraphStar}[1]{\oldparagraph*{#1}\mbox{}}
  \newcommand{\xxxParagraphNoStar}[1]{\oldparagraph{#1}\mbox{}}
\fi
\ifx\subparagraph\undefined\else
  \let\oldsubparagraph\subparagraph
  \renewcommand{\subparagraph}{
    \@ifstar
      \xxxSubParagraphStar
      \xxxSubParagraphNoStar
  }
  \newcommand{\xxxSubParagraphStar}[1]{\oldsubparagraph*{#1}\mbox{}}
  \newcommand{\xxxSubParagraphNoStar}[1]{\oldsubparagraph{#1}\mbox{}}
\fi
\makeatother


\usepackage{longtable,booktabs,array}
\newcounter{none} % for unnumbered tables
\usepackage{calc} % for calculating minipage widths
% Correct order of tables after \paragraph or \subparagraph
\usepackage{etoolbox}
\makeatletter
\patchcmd\longtable{\par}{\if@noskipsec\mbox{}\fi\par}{}{}
\makeatother
% Allow footnotes in longtable head/foot
\IfFileExists{footnotehyper.sty}{\usepackage{footnotehyper}}{\usepackage{footnote}}
\makesavenoteenv{longtable}
\usepackage{graphicx}
\makeatletter
\newsavebox\pandoc@box
\newcommand*\pandocbounded[1]{% scales image to fit in text height/width
  \sbox\pandoc@box{#1}%
  \Gscale@div\@tempa{\textheight}{\dimexpr\ht\pandoc@box+\dp\pandoc@box\relax}%
  \Gscale@div\@tempb{\linewidth}{\wd\pandoc@box}%
  \ifdim\@tempb\p@<\@tempa\p@\let\@tempa\@tempb\fi% select the smaller of both
  \ifdim\@tempa\p@<\p@\scalebox{\@tempa}{\usebox\pandoc@box}%
  \else\usebox{\pandoc@box}%
  \fi%
}
% Set default figure placement to htbp
\def\fps@figure{htbp}
\makeatother



\ifLuaTeX
\usepackage[bidi=basic,shorthands=off]{babel}
\else
\usepackage[bidi=default,shorthands=off]{babel}
\fi
\ifPDFTeX
\else
\babelfont{rm}[]{TeX Gyre Pagella}
\fi
\ifLuaTeX
  \usepackage{selnolig} % disable illegal ligatures
\fi


\setlength{\emergencystretch}{3em} % prevent overfull lines

\providecommand{\tightlist}{%
  \setlength{\itemsep}{0pt}\setlength{\parskip}{0pt}}



 


\usepackage{xcolor}
% ----- COULEUR D'ACCENT BORDEAUX (codes Sciences Po / ESCP) -----
% Pour personnaliser : changer la ligne ci-dessous
% Bleu marine : \definecolor{accent}{HTML}{1F3A5F}
% Bordeaux    : \definecolor{accent}{HTML}{7A1F2B}
% Noir pur    : \definecolor{accent}{HTML}{000000}
% Vert forêt  : \definecolor{accent}{HTML}{1F4D3A}
\definecolor{accent}{HTML}{7A1F2B}
\definecolor{rulegray}{HTML}{B0B0B0}
\definecolor{softgray}{HTML}{555555}

\usepackage{titlesec}
\usepackage{enumitem}
\usepackage{tabularx}
\usepackage{longtable}
\usepackage{array}
\usepackage{paracol}
\usepackage{microtype}
\usepackage{setspace}
\usepackage{ragged2e}

% ----- STYLE DES TITRES DE SECTION -----
% Small caps colorées + filet horizontal bordeaux
% (alternative robuste à \MakeUppercase qui buggue avec accents Unicode)
\titleformat{\section}
  {\sffamily\Large\bfseries\color{accent}\scshape}
  {}{0pt}
  {}
  [\vspace{-4pt}{\color{accent}\titlerule[1pt]}\vspace{2pt}]
\titlespacing*{\section}{0pt}{14pt}{8pt}

\titleformat{\subsection}
  {\sffamily\normalsize\bfseries\color{accent!85!black}}
  {}{0pt}{}
\titlespacing*{\subsection}{0pt}{10pt}{4pt}

% ----- INTERLIGNE LÉGÈREMENT RESSERRÉ -----
\setstretch{1.05}

% ----- PUCES : TIRETS CADRATIN -----
\setlist[itemize]{leftmargin=1.2em, itemsep=2pt, topsep=3pt, label={\color{accent}\textemdash}}
\setlist[enumerate]{leftmargin=1.6em, itemsep=2pt, topsep=3pt}

% ----- PAS DE NUMÉROS DE PAGE EN PAGE 1, DISCRETS APRÈS -----
\usepackage{fancyhdr}
\pagestyle{fancy}
\fancyhf{}
\renewcommand{\headrulewidth}{0pt}
\fancyfoot[R]{\small\color{softgray}\thepage}
\fancypagestyle{plain}{\fancyhf{}\renewcommand{\headrulewidth}{0pt}\fancyfoot[R]{\small\color{softgray}\thepage}}

% ----- COMMANDES PERSONNALISÉES -----
% Ligne de poste avec date alignée à droite
\newcommand{\posteitem}[2]{\noindent\textbf{#1}\hfill{\color{softgray}\small #2}\par}
% Étiquette de classement [CNU X · CNRS Y]
\newcommand{\rank}[1]{{\small\color{accent}\textbf{[#1]}}}
% Item de publication compact
\setlength{\parindent}{0pt}
\setlength{\parskip}{2pt}
\KOMAoption{captions}{tableheading}
\makeatletter
\@ifpackageloaded{caption}{}{\usepackage{caption}}
\AtBeginDocument{%
\ifdefined\contentsname
  \renewcommand*\contentsname{Table des matières}
\else
  \newcommand\contentsname{Table des matières}
\fi
\ifdefined\listfigurename
  \renewcommand*\listfigurename{Liste des Figures}
\else
  \newcommand\listfigurename{Liste des Figures}
\fi
\ifdefined\listtablename
  \renewcommand*\listtablename{Liste des Tables}
\else
  \newcommand\listtablename{Liste des Tables}
\fi
\ifdefined\figurename
  \renewcommand*\figurename{Figure}
\else
  \newcommand\figurename{Figure}
\fi
\ifdefined\tablename
  \renewcommand*\tablename{Table}
\else
  \newcommand\tablename{Table}
\fi
}
\@ifpackageloaded{float}{}{\usepackage{float}}
\floatstyle{ruled}
\@ifundefined{c@chapter}{\newfloat{codelisting}{h}{lop}}{\newfloat{codelisting}{h}{lop}[chapter]}
\floatname{codelisting}{Listing}
\newcommand*\listoflistings{\listof{codelisting}{Liste des Listings}}
\makeatother
\makeatletter
\makeatother
\makeatletter
\@ifpackageloaded{caption}{}{\usepackage{caption}}
\@ifpackageloaded{subcaption}{}{\usepackage{subcaption}}
\makeatother
\usepackage{bookmark}
\IfFileExists{xurl.sty}{\usepackage{xurl}}{} % add URL line breaks if available
\urlstyle{same}
\hypersetup{
  pdflang={fr},
  colorlinks=true,
  linkcolor={accent},
  filecolor={Maroon},
  citecolor={accent},
  urlcolor={accent},
  pdfcreator={LaTeX via pandoc}}


\author{}
\date{}
\begin{document}


\begin{center}
{\sffamily\fontsize{24pt}{28pt}\selectfont\textbf{\color{accent}Georges E\textsc{loundou}}, PhD}\\[4pt]
{\sffamily\large Économiste — Enseignant-chercheur}\\[2pt]
{\itshape Économie appliquée \;\textbar\; Sciences de gestion \;\textbar\; Entrepreneuriat \& Développement \;\textbar\; Analyse de données}\\[6pt]
{\small\color{softgray} Qualifié aux fonctions de Maître de conférences — Section 5 CNU (2026)}\\[8pt]
{\small Métropole du Grand Paris, France \;\textbar\; +33 (0)7 65 28 74 43 \;\textbar\; \href{mailto:eloundou29@gmail.com}{eloundou29@gmail.com}}\\[3pt]
{\small \href{https://orcid.org/0000-0001-8187-5615}{ORCID} \;\textbar\; \href{https://scholar.google.com/citations?user=i_2zSosAAAAJ}{Google Scholar} \;\textbar\; \href{https://www.researchgate.net/profile/Georges-Ngnouwal-Eloundou}{ResearchGate} \;\textbar\; \href{https://www.linkedin.com/in/georges-eloundou}{LinkedIn} \;\textbar\; \href{https://georgeseloundou.github.io}{Site personnel}}
\end{center}

\vspace{4pt}

\section{Profil}\label{profil}

Économiste appliqué positionné à l'intersection de l'économie du
développement, de l'entrepreneuriat, de l'analyse de données et des
sciences de gestion. Mes travaux mobilisent des méthodes mixtes
(économétrie, qualitatif, données primaires) pour analyser les
interactions entre technologies de l'information, énergie, finance et
bien-être dans les économies émergentes. Production scientifique
soutenue de 18 articles dans des revues à comité de lecture (8 CNU B, 6
CNU C · 3 -- ABS 1, 2 ABS 2, 1 ABS 3), ancrage institutionnel
franco-camerounais (Sorbonne Paris Nord, CEPN, CEREG) et activité
éditoriale régulière pour 12 revues internationales ACL.

\section{Positions récentes}\label{positions-ruxe9centes}

\posteitem{Enseignant-chercheur (ATER), IUT de Villetaneuse — Université Sorbonne Paris Nord}{2025 — 2026}
\posteitem{Vacataire (CDD), IUT de Villetaneuse — Université Sorbonne Paris Nord}{2024 — 2025}
\posteitem{Responsable de projet, Centre d'Études et de Recherche sur les Analyses et Politiques Économiques}{2024 — présent}
\posteitem{Chercheur associé, Society for Inclusive and Collaborative Entrepreneurship}{2024 — présent}
\posteitem{Chercheur associé, CEREG — Université de Yaoundé II (membre du comité de rédaction)}{2022 — présent}

\section{Formation}\label{formation}

\hspace{1em}\posteitem{Doctorat (PhD) en Sciences économiques - Analyse des politiques publiques, Économie mathématique}{2024}

\textit{Université de Dschang, CERME - mention Très Honorable}

\hspace{1em}\posteitem{Master II — Monnaie, Banque, Finance et Assurance}{2019}

\textit{Université de Rennes I, France}

\hspace{1em}\posteitem{Master II — Ingénierie économique et financière}{2018}

\textit{Université de Yaoundé II, Cameroun}

\hspace{1em}\posteitem{Licence — Sciences économiques et de gestion}{2016}

\textit{Université de Yaoundé II, Cameroun}

\section{Distinctions \&
qualifications}\label{distinctions-qualifications}

\begin{itemize}
\tightlist
\item
  \textbf{Qualification CNU Section 5 --- Sciences économiques} (2026)
\item
  \textbf{Mention Très Honorable} au Doctorat (2024)
\item
  \textbf{Lauréat de la Bourse Erasmus Mundus} (2018)
\end{itemize}

\newpage

\section{Production scientifique}\label{production-scientifique}

\vspace{-2pt}

\begin{center}
\small
\begin{tabular}{l c l}
\toprule
\textbf{Type de production} & \textbf{Total} & \textbf{Détail} \\
\midrule
Articles dans revues à comité de lecture & 18 & 8 CNU B · 6 CNU C · 3 NC \\
Working papers & 7 & dont 1 publié \\
Chapitre d'ouvrage & 1 & Éditions Cheikh Anta Diop (2023) \\
Rapports de recherche & 6 & PNDP, OCDE-UA, RSE \\
Communications colloques internationaux & 6 & France, Cameroun, Maroc \\
Séminaires, journées d'étude, workshops & 9 & — \\
\bottomrule
\end{tabular}
\end{center}

\vspace{-2pt}

\noindent\textit{Rayonnement} : plus de 8,000 lectures et 100 citations
recensées sur ResearchGate.

\subsection{Articles dans des revues à comité de
lecture}\label{articles-dans-des-revues-uxe0-comituxe9-de-lecture}

\begin{enumerate}[leftmargin=1.6em, itemsep=3pt]

\item Ngnouwal Eloundou, G. \textit{et al.} (2026). Can we build inclusion in Africa? The key role of infrastructure'' \textit{Economics Bulletin}. \rank{CNU B · CNRS 3} \href{http://www.accessecon.com/Pubs/EB/2026/Volume46/EB-26-V46-I1-P1.pdf}{Lien}

\item Ngnouwal Eloundou, G. \textit{et al.} (2026). Linking country stability and informality: empirical evidence from developing countries. \textit{Oxford Development Studies}. \rank{CNU B · CNRS 3 · ABS 2} \href{https://doi.org/10.1080/13600818.2026.2654469}{DOI}

\item Ngnouwal Eloundou, G. \textit{et al.} (2025). The intensive role of FDI on women's empowerment in developing countries. \textit{International Journal of Finance \& Economics}. \rank{CNU B · CNRS 3 · ABS 3} \href{https://doi.org/10.1002/ijfe.70097}{DOI}

\item Ngnouwal Eloundou, G. \textit{et al.} (2025). Good or bad? Environmental policy and women's political empowerment in developing countries. \textit{Economics Bulletin}. \rank{CNU B · CNRS 3} \href{https://econpapers.repec.org/article/eblecbull/eb-23-00209.htm}{Lien}

\item Ngnouwal Eloundou, G. \textit{et al.} (2025). Does social media foster informality in developing countries? \textit{European Review of Service Economics and Management}. \rank{CNU C · CNRS 4} \href{https://doi.org/10.48611/isbn.978-2-406-18573-4.p.0071}{DOI}

\item Ngnouwal Eloundou, G. \textit{et al.} (2025). Impact of nation branding on economic growth in Sub-Saharan Africa. \textit{The Bottom Line}. \href{https://doi.org/10.1108/BL-07-2024-0088}{DOI}

\item Ngnouwal Eloundou, G. \textit{et al.} (2025). On the path to accelerating industrialization in Africa: the role of fiscal decentralization. \textit{Journal of the Knowledge Economy}. \rank{CNU C · CNRS 4 · ABS 1} \href{https://doi.org/10.1007/s13132-025-02602-2}{DOI}

\item Ngnouwal Eloundou, G. \textit{et al.} (2025). Energy poverty and entrepreneurship: evidence from Sub-Saharan Africa. \textit{The Bottom Line}. \href{https://doi.org/10.1108/BL-09-2024-0169}{DOI}

\item Ngnouwal Eloundou, G. \textit{et al.} (2024). Black shadow: effect of energy poverty on informal economy in developing countries. \textit{The Journal of Energy and Development}. \rank{CNU B · CNRS 3} \href{https://doi.org/10.56476/jed.v49i1.21}{DOI}

\item Ngnouwal Eloundou, G. \textit{et al.} (2024). Does social media contribute to economic growth? \textit{Journal of the Knowledge Economy}, 1–41. \rank{CNU C · CNRS 4 · ABS 1} \href{https://doi.org/10.1007/s13132-023-01419-1}{DOI}

\item Ngnouwal Eloundou, G. \textit{et al.} (2024). Alcohol consumption in developing countries: does information and communication technology (ICT) diffusion matter? \textit{Journal of International Development}. \rank{CNU B · CNRS 3 · ABS 2} \href{https://doi.org/10.1002/jid.3858}{DOI}

\item Ngnouwal Eloundou, G. \textit{et al.} (2023). In pursuit of happiness in developing countries: does the diffusion of ICT matter? \textit{Economics Bulletin}, 43(4), 1730–1740. \rank{CNU B · CNRS 3} \href{https://econpapers.repec.org/article/eblecbull/eb22-00799.htm}{Lien}

\item Ngnouwal Eloundou, G. \textit{et al.} (2023). Foreign direct investment resilience during times of crisis: a comparative analysis between selected eastern European and African countries. \textit{International Politics}. \href{https://doi.org/10.18485/iipe_mp.2023.74.1189.4}{DOI}

\item Ngnouwal Eloundou, G. \textit{et al.} (2023). Does social media drive remittances in Africa? \textit{African Development Review}. \rank{CNU C · CNRS 4} \href{https://doi.org/10.1111/1467-8268.12717}{DOI}

\item Ngnouwal Eloundou, G. \textit{et al.} (2023). Does social media improve women's political empowerment in Africa? \textit{Telecommunications Policy}, 47(9), 102624. \rank{CNU C · CNRS 4 · ABS 1} \href{https://doi.org/10.1016/j.telpol.2023.102624}{DOI}

\item Ngnouwal Eloundou, G. \textit{et al.} (2023). Is access to electricity a relevant determinant of happiness in developing countries? \textit{The Journal of Energy and Development}. \rank{CNU B · CNRS 3} \href{https://www.jstor.org/stable/27284853}{Lien}

\item Ngnouwal Eloundou, G. \textit{et al.} (2022). Impact of access to electricity on internal conflicts in Africa: does income inequality matter? \textit{African Development Review}, 34(3), 395–409. \rank{CNU C · CNRS 4} \href{https://doi.org/10.1111/1467-8268.12664}{DOI}

\end{enumerate}

\subsection{Chapitre d'ouvrage}\label{chapitre-douvrage}

\begin{itemize}
\tightlist
\item
  Ngnouwal Eloundou, G. \textit{et al.} (2024). Quels sont les
  déterminants pertinents des fragilités des pays en développement ? In
  B. E. Ongo Nkoa (Coord.),
  \textit{Le développement économique à la croisée des réflexions : illustrations par des données africaines}.
  Éditions Cheikh Anta Diop (ouvrage collectif, 46 contributeurs, 700
  pages).
\end{itemize}

\subsection{Working paper publié}\label{working-paper-publiuxe9}

\begin{itemize}
\tightlist
\item
  Ngnouwal Eloundou, G. \textit{et al.} (2022). Does infrastructure
  development contribute to poverty reduction in Sub-Saharan Africa?
  \textit{SSRN Working Paper}.
  \href{https://ssrn.com/abstract=4397985}{SSRN}
\end{itemize}

\subsection{Publications en cours (soumissions \&
révisions)}\label{publications-en-cours-soumissions-ruxe9visions}

\begin{enumerate}[leftmargin=1.6em, itemsep=2pt]
\item What makes women entrepreneurs in the informal sector resilient?. International Journal of Entrepreneurial Behaviour and Research · ABS 3
\item Intellectual Property and Sustainable Development. Review of Development Economics · ABS 3
\item Institution, natural resources and inequality. Resources Policy · ABS 2
\item Do social media improve entrepreneurship in developing economies? The Journal of Entrepreneurship · ABS 1
\item Effect of entrepreneurship on irregular migration: does regulation matter?
\item Influence of entrepreneurial perceptions on entrepreneurial activity
\end{enumerate}

\subsection{Rapports de recherche \&
expertise}\label{rapports-de-recherche-expertise}

\begin{itemize}
\tightlist
\item
  \textbf{République du Cameroun (Juin 2026)} --- Coauteur rapport «
  Rapport d'étude pour l'évaluation des impacts socio-économiques des
  changements climatiques sur les populations vivant en périphérie du
  Nyong » (coord. Désiré Avom).
\item
  \textbf{OCDE \& Commission de l'Union Africaine (2024)} --- Coauteur
  du chapitre 3 « Compétences sur l'exploitation minière en Afrique
  centrale » du rapport
  \textit{Dynamiques du développement durable en Afrique 2024 : compétences, emplois et productivité}
  (coord. B. E. Ongo Nkoa).
\item
  \textbf{PNDP -- Citizen Report Cards (2023)} --- Coauteur des rapports
  d'évaluation des services publics (eau, santé, éducation, services
  communaux) pour les communes de Njimom, Santchou, Fongo-Tongo et
  Penka-Michel.
\item
  \textbf{Cabinet d'Étude Multipolaire (2023)} --- Coauteur du rapport «
  Perception des riverains des projets publics de la RSE des entreprises
  étrangères » (sites de Batchenga et Ntui).
\end{itemize}

\subsection{Communications dans des colloques
internationaux}\label{communications-dans-des-colloques-internationaux}

\begin{itemize}[itemsep=2pt]
\item \textbf{8\textsuperscript{th} Entrepreneurship Innovation and Governance Conference (ENIG)} — EDC Paris Business School, 11 juin 2026.
\item \textbf{6\textsuperscript{th} International Conference on Development Economics (ICDE)} — Université Rouen Normandie, 2–3 juillet 2026.
\item \textbf{45\textsuperscript{es} Journées de l'Association d'Économie Sociale} — CEPN UR 7234, Sorbonne Paris Nord, 3–4 septembre 2026.
\item \textbf{14\textsuperscript{e} Congrès de l'Académie de l'Entrepreneuriat et de l'Innovation (AEI 2025)} : « L'entrepreneuriat face aux transitions » — Aix-en-Provence, juin 2025.
\item \textbf{5\textsuperscript{e} Colloque International sur l'Écosystème Entrepreneurial} : « Entreprendre aujourd'hui : profils d'entrepreneurs et carrières entrepreneuriales » — Douala, avril 2025.
\item \textbf{4\textsuperscript{e} Conférence Internationale sur la Francophonie Économique} : « L'avenir des PME francophones sur les marchés mondiaux » — Yaoundé, mars 2024.
\item \textbf{4\textsuperscript{e} Colloque International sur l'Écosystème Entrepreneurial} : « Cultiver la résilience entrepreneuriale dans un environnement changeant » — Casablanca, avril 2024.
\end{itemize}

\section{Activités éditoriales \& rayonnement
scientifique}\label{activituxe9s-uxe9ditoriales-rayonnement-scientifique}

\textbf{Reviewing régulier} pour :
\textit{Journal of International Development · ABS 2} ·
\textit{Journal of the Knowledge Economy · ABS 1} ·
\textit{African Development Review} · \textit{The Bottom Line} ·
\textit{Applied Economics Letters · ABS 1} ·
\textit{World Development Perspectives} · \textit{Economics Bulletin} ·
\textit{Politics and Policy · ABS 3} ·
\textit{South African Journal of Economics · ABS 1} ·
\textit{Review of Development Economics · ABS 2} ·
\textit{Journal of Entrepreneurship in Emerging Economies · ABS 1} ·
\textit{ Australian Economic Papers · ABS 1}.

\textbf{Organisation scientifique} --- Co-organisateur et porteur de la
session thématique « Entrepreneuriat informel, inclusif et durable :
vers une reconfiguration des pratiques et des théories de
l'Entre-Prendre » lors du
\textbf{18\textsuperscript{e} CIFEPME, Bruxelles 2026}.

\textbf{Affiliations à des laboratoires \& centres de recherche} :

\begin{itemize}[itemsep=1pt]
\item \textbf{France} — CEPN (Centre d'Économie de Paris Nord, UR 7234) · AEI (Académie de l'Entrepreneuriat et de l'Innovation) · AFEDEV (Association Française des Economistes du Développement).
\item \textbf{Cameroun} — CEREG (Université de Yaoundé II) · LAREFA (Université de Dschang) · CERAPE · ACEDA (Association des chercheurs en économie du développement appliquée).
\end{itemize}

\section{Enseignement}\label{enseignement}

Plus de \textbf{400 heures équivalent TD enseignées} depuis 2021, sur
des publics variés (Licence, BUT, Bachelor finance) en France et au
Cameroun. Pédagogie active fondée sur l'étude de cas, l'accompagnement
de projets entrepreneuriaux et le tutorat de soutenances SAE. Cours
actuels alignés sur les codes business school : entrepreneuriat,
management, accompagnement du changement, business plan, finance de
marché.

\vspace{2pt}

\begin{center}
\footnotesize
\renewcommand{\arraystretch}{1.25}
\begin{tabularx}{\textwidth}{@{}l l X r r@{}}
\toprule
\textbf{Année} & \textbf{Niveau} & \textbf{Cours} & \textbf{Vol. (h)} & \textbf{Eff.} \\
\midrule
\multicolumn{5}{l}{\textit{IUT de Villetaneuse — Université Sorbonne Paris Nord, Département Management}} \\
\midrule
2025–2026 & BUT3 GEII & Entrepreneuriat (4 groupes) & 60 & 103 \\
2025–2026 & BUT3 GEMA & Accompagnement du changement & 20 & 38 \\
2025–2026 & BUT3 GEMA & Management des activités / relations commerciales & 50 & 38 \\
2025–2026 & BUT3 GMIE & Management & 12 & 15 \\
2025–2026 & BUT3 GEMA & Projet personnel et professionnel & 2 & 23 \\
2025–2026 & BUT2 GEMA App. & Outils et pilotage : Business Plan & 14 & 10 \\
2025–2026 & BUT1 GEII & Gestion de projet & 25 & 28 \\
2025–2026 & BUT1 GEMA & Jury de soutenance SAE & 4 & — \\
2024–2025 & BUT3 GEMA & Accompagnement du changement & 10 & 12 \\
\midrule
\multicolumn{5}{l}{\textit{Yaoundé School of Economics and Management — Département Finance}} \\
\midrule
2022–2024 & Bachelor 3 & Évaluation des actifs financiers & 30 & 15 \\
2022–2024 & Bachelor 3 & Gestion de portefeuille & 30 & 15 \\
\midrule
\multicolumn{5}{l}{\textit{Institut Universitaire des Sciences et Techniques de Yaoundé}} \\
\midrule
2022–2023 & Licence 1 & Économie générale & 30 & 200 \\
\midrule
\multicolumn{5}{l}{\textit{Monitorat travaux dirigés — Universités de Dschang et de Yaoundé II}} \\
\midrule
2018–2023 & TD Licence & Macroéconomie, microéconomie, statistiques, économie industrielle & — & — \\
\bottomrule
\end{tabularx}
\end{center}

\section{Expertise \& projets de
terrain}\label{expertise-projets-de-terrain}

\begin{itemize}
\tightlist
\item
  \textbf{PNUD --- Évaluation à mi-parcours des ODD en Afrique
  centrale.} Membre de l'équipe de rédaction du rapport « Opportunités
  et contraintes vers une intégration réussie ».
\item
  \textbf{PNDP -- Projet SCORECARD2, Ouest Cameroun.} Coordonnateur
  adjoint de la collecte de données et coauteur des rapports
  d'évaluation de l'offre de services de base (eau, santé, éducation)
  --- communes de Njimom, Penka-Michel, Fongo-Tongo, Massangam.
\item
  \textbf{Cabinet d'Étude Multipolaire --- Évaluation de la RSE
  (Batchenga, Ntui).} Co-concepteur des outils de collecte qualitative
  et quantitative, analyste, coauteur du rapport.
\item
  \textbf{S4ICE --- Entrepreneuriat culturel en Afrique centrale
  (Cameroun, Tchad).} Coordonnateur adjoint, concepteur des outils
  d'enquête, analyste mixte, coauteur.
\end{itemize}

\section{Compétences}\label{compuxe9tences}

\columnratio{0.5}
\begin{paracol}{2}

\textbf{Méthodes \& logiciels d'analyse}\\
R · Stata · Python · SPSS · NVivo · SQL · Excel\\

\textbf{Collecte de données primaires}\\
ODK · CSPro · Google Forms · conception de questionnaires et guides d'entretien · animation de focus groups\\

\textbf{Rédaction \& présentation}\\
\LaTeX{} · Beamer · Pack MS Office · Moodle · Quarto · HTML

\switchcolumn

\textbf{Gestion de projets}\\
MS Project · GanttProject · planification stratégique · suivi-évaluation · mobilisation de fonds\\

\textbf{Conseil \& expertise}\\
Stratégies de développement durable · responsabilité sociétale des entreprises (RSE)\\

\textbf{Langues}\\
Français (langue maternelle) · Anglais (niveau intermédiaire avancé)

\end{paracol}




\end{document}
